\section*{Hoofdstuk \ref{ch:g11:deriv} --- Afgeleiden: een eerste kennismaking}

\begin{solution}{exo:g11:deriv:1}
Voor $f(x) = x^2$ in $a = 2$:
\[
\frac{(2+h)^2 - 4}{h} = \frac{4h + h^2}{h} = 4 + h
\xrightarrow[h\to0]{} 4,
\]
dus $f'(2) = 4$. Voor $g(x) = x^2 + 3x$ in $a = 1$: $g(1) = 4$ en
\[
\frac{(1+h)^2 + 3(1+h) - 4}{h} = \frac{5h + h^2}{h} = 5 + h
\xrightarrow[h\to0]{} 5,
\]
dus $g'(1) = 5$.
\end{solution}

\begin{solution}{exo:g11:deriv:2}
$f'(x) = 12x^2 - 4x + 1$.

$g'(x) = \frac{5}{2\sqrt x} - \frac{3}{x^2}$ (op $\intoo{0}{+\infty}$).

$h'(x) = 2(x^2 - 3) + (2x+1)(2x) = 2x^2 - 6 + 4x^2 + 2x = 6x^2 + 2x - 6$
(productregel).
\end{solution}

\begin{solution}{exo:g11:deriv:3}
$f(x) = x^2 - 3x$, $a = 2$: $f(2) = -2$ en $f'(x) = 2x - 3$, dus
$f'(2) = 1$; \omterm{def:g11:deriv:tangent}{raaklijn} $y = -2 + (x - 2) = x - 4$.

$f(x) = \frac1x$, $a = 1$: $f(1) = 1$, $f'(1) = -1$; \omterm{def:g11:deriv:tangent}{raaklijn}
$y = 1 - (x - 1) = 2 - x$.

$f(x) = \sqrt x$, $a = 4$: $f(4) = 2$, $f'(4) = \frac{1}{2\sqrt4} =
\frac14$; \omterm{def:g11:deriv:tangent}{raaklijn} $y = 2 + \frac14(x - 4) = \frac{x}{4} + 1$.
\end{solution}

\begin{solution}{exo:g11:deriv:4}
$f'(x) = \frac{(x-1) - (x+2)}{(x-1)^2} = \frac{-3}{(x-1)^2}$.

$g'(x) = \frac{2x(x^2+1) - x^2 \times 2x}{(x^2+1)^2}
= \frac{2x}{(x^2+1)^2}$.

$h'(x) = -\frac{2x+1}{(x^2+x+1)^2}$ (quotiëntregel met $u = 1$).
\end{solution}

\begin{solution}{exo:g11:deriv:5}
$f'(x) = 3x^2 - 12x + 9 = 3(x^2 - 4x + 3) = 3(x-1)(x-3)$: positief buiten
$\intcc{1}{3}$, negatief erbinnen. Dus is $f$ \omterm{def:g11:func:monotone}{stijgend} op
$\intoc{-\infty}{1}$, \omterm{def:g10:functions:variations}{dalend} op $\intcc{1}{3}$ en \omterm{def:g11:func:monotone}{stijgend} op
$\intco{3}{+\infty}$, met een lokaal \omterm{def:g10:functions:extrema}{maximum} $f(1) = 1 - 6 + 9 - 2 = 2$ en
een lokaal \omterm{def:g10:functions:extrema}{minimum} $f(3) = 27 - 54 + 27 - 2 = -2$.
\end{solution}

\begin{solution}{exo:g11:deriv:6}
Op $\intoo{-1}{+\infty}$:
\[
f'(x) = \frac{2x(x+1) - (x^2+3)}{(x+1)^2}
= \frac{x^2 + 2x - 3}{(x+1)^2}
= \frac{(x+3)(x-1)}{(x+1)^2}.
\]
Voor $x > -1$ zijn zowel $x + 3$ als $(x+1)^2$ positief, dus heeft $f'(x)$
het teken van $x - 1$: $f$ is \omterm{def:g10:functions:variations}{dalend} op $\intoc{-1}{1}$ en \omterm{def:g11:func:monotone}{stijgend} op
$\intco{1}{+\infty}$, met \omterm{def:g10:functions:extrema}{minimum} $f(1) = \frac{4}{2} = 2$.
\end{solution}

\begin{solution}{exo:g11:deriv:7}
$f'(x) = 2x - 4$.

\emph{1.} Horizontale \omterm{def:g11:deriv:tangent}{raaklijn}: $f'(x) = 0$ in $x = 2$; het punt is
$(2, f(2)) = (2, 1)$ (de top van de \omterm{def:g11:quad:parabola}{parabool}).

\emph{2.} Evenwijdig met $y = 2x + 1$ betekent \omterm{def:g10:reffunc:affine}{richtingscoëfficiënt} $2$:
$2x - 4 = 2$, dus $x = 3$, wat het punt $(3, f(3)) = (3, 2)$ geeft.
\end{solution}

\begin{solution}{exo:g11:deriv:8}
Met breedte $x$ (twee omheinde zijden) en lengte $60 - 2x$ is
$A(x) = x(60 - 2x) = 60x - 2x^2$ op $\intoo{0}{30}$. Dan is
$A'(x) = 60 - 4x$, positief vóór $x = 15$ en negatief erna: de oppervlakte
is maximaal voor $x = 15$, wat een weide van $15 \times 30$ m met
oppervlakte $450$ m$^2$ geeft. De topformule
$\alpha = -\frac{60}{2 \times (-2)} = 15$ uit \cref{ch:g11:quad} geeft
hetzelfde antwoord zonder \omterm{def:g11:deriv:derivative}{afgeleiden}.
\end{solution}

\begin{solution}{exo:g11:deriv:9}
Zij $f(x) = \frac{x+1}{2} - \sqrt x$ op $\intoo{0}{+\infty}$. Dan is
\[
f'(x) = \frac12 - \frac{1}{2\sqrt x} = \frac{\sqrt x - 1}{2\sqrt x},
\]
wat het teken heeft van $\sqrt x - 1$: negatief op $\intoo{0}{1}$, positief
op $\intoo{1}{+\infty}$. Dus daalt $f$ eerst en stijgt ze daarna, met
\omterm{def:g10:functions:extrema}{minimum} $f(1) = 1 - 1 = 0$. Bijgevolg is $f(x) \geq 0$ voor alle $x > 0$,
dus $\sqrt x \leq \frac{x+1}{2}$, met gelijkheid alleen in $x = 1$.
\end{solution}

\begin{solution}{exo:g11:deriv:10}
Uit $\pi r^2 h = 500$ volgt $h = \frac{500}{\pi r^2}$, dus
\[
S(r) = 2\pi r^2 + 2\pi r \cdot \frac{500}{\pi r^2}
= 2\pi r^2 + \frac{1000}{r},
\qquad
S'(r) = 4\pi r - \frac{1000}{r^2}.
\]
$S'(r) = 0$ wanneer $4\pi r^3 = 1000$, dus $r^3 = \frac{250}{\pi}$ en
$r = \sqrt[3]{250/\pi}$ ($\approx 4.3$ cm). Omdat
$S'(r) = \frac{4\pi r^3 - 1000}{r^2}$ vóór die waarde negatief is en erna
positief, gaat het om een \omterm{def:g10:functions:extrema}{minimum}.
\end{solution}

\begin{solution}{exo:g11:deriv:11}
De \omterm{def:g11:deriv:tangent}{raaklijn} aan $y = x^2$ in het punt met \omterm{def:g10:coordgeom:system}{abscis} $a$ is
$y = a^2 + 2a(x - a) = 2ax - a^2$. Ze gaat door $P(1, -3)$ wanneer
$-3 = 2a - a^2$, dus $a^2 - 2a - 3 = 0$, met \omterm{def:g11:quad:discriminant}{wortels} $a = 3$ en $a = -1$
($\Delta = 16$). Er zijn dus \emph{twee} \omterm{def:g11:deriv:tangent}{raaklijnen} door $P$:
\[
y = 6x - 9 \quad (a = 3),
\qquad
y = -2x - 1 \quad (a = -1).
\]
\end{solution}

\begin{solution}{exo:g11:deriv:12}
Uit \cref{ex:g11:deriv:study}: $f$ stijgt tot het lokale \omterm{def:g10:functions:extrema}{maximum}
$f(-1) = 3$, daalt tot het lokale \omterm{def:g10:functions:extrema}{minimum} $f(1) = -1$, en stijgt daarna
opnieuw; voor grote positieve (respectievelijk negatieve) $x$ overheerst
$x^3$ en neemt $f$ willekeurig grote (respectievelijk kleine) waarden aan.
Trek je horizontale rechten $y = k$ over de \omterm{def:g10:functions:table}{verlooptabel}, dan geldt:
\begin{itemize}
  \item $k > 3$ of $k < -1$: de rechte ontmoet de kromme één keer ---
        $1$ oplossing;
  \item $k = 3$ of $k = -1$: de rechte raakt een keerpunt en kruist één
        andere tak --- $2$ oplossingen;
  \item $-1 < k < 3$: de rechte kruist alle drie de takken ---
        $3$ oplossingen.
\end{itemize}
\end{solution}

\begin{solution}{pb:g11:deriv:1}
\textbf{1.} $f'(x) = 6x - 5$; $g'(x) = 3x^2 - 12$. Uit de definitie:
$\frac{(1 + h)^2 - 1}{h} = 2 + h \to 2$: de \omterm{def:g11:deriv:derivative}{afgeleide} van $x^2$ in $1$ is
$2$.

\textbf{2.} $g(2) = 8 - 24 = -16$ en $g'(2) = 12 - 12 = 0$: de \omterm{def:g11:deriv:tangent}{raaklijn} is
de horizontale rechte $y = -16$. Een horizontale \omterm{def:g11:deriv:tangent}{raaklijn} signaleert een
kritiek punt --- hier een lokaal \omterm{def:g10:functions:extrema}{minimum}, zoals de \omterm{def:g10:functions:table}{verlooptabel} van vraag 5
bevestigt.

\textbf{3.} \omterm{def:g10:reffunc:affine}{Richtingscoëfficiënt} $2a$ in $(a, a^2)$:
$y = a^2 + 2a(x - a) = 2ax - a^2$.

\textbf{4.} In $x = 0$: $y = -a^2$: de \omterm{def:g11:deriv:tangent}{raaklijn} snijdt de as op diepte $a^2$
onder de oorsprong --- precies even ver \emph{onder} nul als $P$ erboven
ligt. Recept: om de \omterm{def:g11:deriv:tangent}{raaklijn} te tekenen in een punt op hoogte $h$, verbind
je het met het aspunt op diepte $h$ onder de oorsprong. (Aan de tekentafel
heb je geen \omterm{def:g11:deriv:derivative}{afgeleide} nodig.)

\textbf{5.} $g'(x) = 3(x^2 - 4)$: positief op $\intoo{-\infty}{-2}$,
negatief op $\intoo{-2}{2}$, positief daarvoorbij. $g$ stijgt tot een lokaal
\omterm{def:g10:functions:extrema}{maximum} $g(-2) = 16$, zakt tot een lokaal \omterm{def:g10:functions:extrema}{minimum} $g(2) = -16$ en stijgt
daarna.

\textbf{6.} $F\left(0, \frac14\right)$ en $T(0, -a^2)$:
$FT = \frac14 + a^2$.

\textbf{7.} $FP = a^2 + \frac14 = FT$: gelijkbenig in $F$, dus zijn de
basishoeken gelijk: $\widehat{FTP} = \widehat{FPT}$.

\textbf{8.} De invallende straal in $P$ is verticaal, en dus evenwijdig met
de rechte $(TF)$ (de $y$-as). De \omterm{def:g11:deriv:tangent}{raaklijn} $(TP)$ kruist beide evenwijdige
rechten, zodat de hoek tussen de straal en de \omterm{def:g11:deriv:tangent}{raaklijn} gelijk is aan
$\widehat{FTP}$ (verwisselende hoeken) --- en die is gelijk aan
$\widehat{FPT}$, de hoek tussen $[PF]$ en de \omterm{def:g11:deriv:tangent}{raaklijn}. Gelijke hoeken aan
weerszijden van de \omterm{def:g11:deriv:tangent}{raaklijn} in $P$.

\textbf{9.} Volgens de weerkaatsingswet vertrekt de straal die verticaal in
$P$ aankomt langs de rechte die aan de andere kant dezelfde hoek maakt ---
volgens vraag 8 dus langs $[PF]$: door het \omterm{pb:g11:quad:1}{brandpunt}, wat $a$ ook is.
Omgekeerd verlaat licht dat in $F$ ontstaat elk punt van de spiegel
evenwijdig met de as. Schotels vangen op, koplampen stralen uit: bewezen.

\textbf{10.} $a = 1$: \omterm{def:g11:deriv:tangent}{raaklijn} $y = 2x - 1$, $T(0, -1)$;
$FP = \sqrt{1 + \left(1 - \frac14\right)^2} = \sqrt{\frac{25}{16}} =
\frac54$; $FT = \frac14 + 1 = \frac54$: gelijkbenig, zoals beloofd.

\textbf{11.} De \omterm{def:g11:deriv:tangent}{raaklijn} in $x_n$ is $y = f(x_n) + f'(x_n)(x - x_n)$; ze
snijdt $y = 0$ waar $x = x_n - \frac{f(x_n)}{f'(x_n)}$.

\textbf{12.} $x_1 = 1 - \frac{-1}{2} = \frac32$;
$x_2 = \frac32 - \frac{1/4}{3} = \frac{17}{12}$; $x_3 = \frac{577}{408}$:
de \omterm{def:g10:numbers:approx}{benaderingen} van $\sqrt2$ van Heron, dus de machine met het vaste punt
uit \cref{pb:g10:functions:1}. Identiteit:
$x - \frac{x^2 - 2}{2x} = \frac{2x^2 - x^2 + 2}{2x} = \frac{x^2 + 2}{2x}
= \frac12\left(x + \frac2x\right)$ --- het recept van Heron \emph{is} de
methode van Newton toegepast op $x^2 - 2$, twee millennia te vroeg.

\textbf{13.} $x_1 = 5 - \frac{125 - 100}{75} = \frac{14}{3} \approx 4.6667$;
$x_2 \approx 4.6667 - \frac{1.63}{65.33} \approx 4.6417$. Rekenmachine:
$\sqrt[3]{100} = 4.6416\ldots$ --- vier juiste decimalen in twee stappen.

\textbf{14.} $f'(2) = 0$: de \omterm{def:g11:deriv:tangent}{raaklijn} in het startpunt is horizontaal (vraag
2) en snijdt de as nooit --- de formule deelt door nul. Voorbehoud: start
Newton weg van kritieke punten (en, algemener, dicht genoeg bij de \omterm{def:g11:quad:discriminant}{wortel}
die je zoekt).

\textbf{15.} Bisectie wint per stap één binair cijfer; Newton
\emph{verdubbelt} ruwweg het aantal juiste cijfers per stap (hier
$1 \to 3 \to 6$). Wanneer elke stap duur is --- miljoenen parameters,
navigatie in real time --- verslaat verdubbelen het stapje voor stapje
vorderen, en daarom draait er zoveel silicium op de \omterm{def:g11:deriv:tangent}{raaklijn}.

\textbf{16.} $d^2 = 900 - w^2$, dus $S(w) = w(900 - w^2) = 900w - w^3$, voor
$0 < w < 30$.

\textbf{17.} $S'(w) = 900 - 3w^2 = 0$ in $w = \sqrt{300} = 10\sqrt3
\approx 17.3$ cm ($S' > 0$ ervóór, $< 0$ erna: een \omterm{def:g10:functions:extrema}{maximum}). Dan is
$d = \sqrt{600} = 10\sqrt6 \approx 24.5$ cm.

\textbf{18.} $\frac{d^2}{w^2} = \frac{600}{300} = 2$, dus $d = w\sqrt2$: de
hoogte is de breedte maal $\sqrt2$ --- de verhouding van de timmerman (en de
constructie met de diameter in drieën levert ze met een passer, zonder
algebra op de werf).

\textbf{19.} Vierkante balk: $w = d = \sqrt{450} \approx 21.2$ cm, met
stijfheid $w^3 = 450^{3/2} \approx 9\,546$. Optimale balk:
$S = 10\sqrt3 \times 600 = 6\,000\sqrt3 \approx 10\,392$. Winst: ongeveer
$8.9\,\%$ meer stijfheid uit dezelfde stam --- de \omterm{def:g11:deriv:derivative}{afgeleide} betaalt de
zagerij terug.

\textbf{20.} Meetkunde: de gelijke hoeken van de \omterm{def:g11:deriv:tangent}{raaklijn} maakten van een
paraboloïde een lichtverzamelaar. Numeriek rekenen: langs \omterm{def:g11:deriv:tangent}{raaklijnen}
afglijden lost \omterm{def:g10:algebra:equation}{vergelijkingen} op met verdubbelende cijfers. Optimaliseren:
een verdwijnende \omterm{def:g11:deriv:derivative}{afgeleide} vond de sterkste rechthoek in de stam. Eén
object, drie stielen --- en de convexiteit van jaar 12 zal daarbovenop
vertellen aan welke kant van elke \omterm{def:g11:deriv:tangent}{raaklijn} de kromme ligt.
\end{solution}
